\chapter{SCF rescue (menu 9)}
\label{ch:menu9}
\menulabel{9}

\section{Purpose}

Read one output file whose SCF struggled, name what happened (triage), and --
if you give the matching input file -- write corrected input files for you,
each with an auditable list of what changed and why.

\section{What you need}

\begin{itemize}
  \item the ORCA output file of the failed or suspicious SCF;
  \item optionally the matching input file: with it, \menu{9} writes corrected
    inputs; without it, you get the triage only.
\end{itemize}

\section{Walkthrough}

\lstinputlisting[style=fbkcode, firstline=54, lastline=69,
  title={The menu-9 example, part 1 (an SCF that did not converge: findings and
  corrected inputs)}]
  {examples/expected/m9_scf.txt}

\lstinputlisting[style=fbkcode, firstline=122, lastline=135,
  title={Menu 9 on a healthy run (part 2): the verbatim convergence block, then
  silence from the triage}]
  {examples/expected/m9_scf.txt}

\section{What you get}

First the \code{SCF CONVERGENCE} block verbatim, exactly as ORCA printed it,
whenever the output carries one -- the informational rows (density changes,
DIIS error) are shown without any judgement, because the check mode decides
which rows are enforced (on a healthy run a DIIS row far above its tolerance
is informational, not a finding). The aborted run of the example prints no
summary block, so it goes straight to the findings. Then findings from the
triage (each with the measured numbers -- cycle counts, DIIS
errors, achieved-versus-tolerance comparisons -- and a suggested action; the
repertoire spans no convergence, a crash without a verdict, pseudo-convergence,
a long run, a DIIS-error rebound after the AO-DIIS switch, an energy trajectory
that oscillates without collapsing, and unmet enforced criteria), then
one corrected input per applicable fix, written next to your input as
\file{name.fix\_\textit{variant}.inp} (three variants exist: the damping
\code{SlowConv} one, the two-step pre-SCF route, and a \code{TRAH} variant for
a rising or oscillating trajectory -- ORCA has no keyword spelled ARH, TRAH is
its trust-region augmented-Hessian second-order SCF with AutoTRAH on by
default). The example wrote two:

\begin{itemize}
  \item \file{scf\_noconv.fix\_slowconv.inp} -- appends the manual's damping
    keyword, with the manual's caution carried in the file as a comment
    (\code{SlowConv} can converge to a solution closer to the initial guess);
  \item \file{scf\_noconv.fix\_prescf.inp} -- the manual's two-step route: a
    cheap pre-SCF whose only purpose is better starting orbitals, then the
    original job with \code{moread} and \code{\%moinp} pointing at the pre-SCF
    orbitals. Your \code{\%scf} settings are preserved (kept as comments in
    step~1 and used unchanged in step~2).
\end{itemize}

\section{How to read the numbers}

\begin{readbox}
Read the triage finding first: it says what pattern was recognised and quotes
the measured numbers from your file. Then read each proposal's rationale and
change list -- nothing is changed silently, and every corrected file states its
own uncertainty (the pre-SCF route ``promises nothing''; the damping route is
the weaker of the two when the starting orbitals are the true problem).
\end{readbox}

\section{Worked example}

\begin{itemize}
  \item The session above is the failure case: an SCF stopped after 2 cycles
    with a large final DIIS error.
  \item The negative control is a healthy run: point \menu{9} at
    \file{n2\_hf\_clean.out} (a clean 8-cycle SCF) and the triage answers
    ``SCF looks healthy: no triage finding.'' -- no noise, no invented
    problems.
\end{itemize}

\section{Caveats, refusals and related sections}

\begin{refusalbox}
\begin{itemize}
  \item \textbf{It will not propose merely raising \code{MaxIter}}: the manual
    states plainly that this does not help in many cases; the program says so
    as a refusal instead of writing that change.
  \item Your original files are never modified; every product is a new file.
  \item The triage reads numbers printed by ORCA; when a number is missing
    from the output (a truncated file, a different print level), the
    corresponding criterion is skipped rather than guessed.
  \item The proposed inputs keep your file's structure (comment lines,
    \code{\%scf} block) and add only what the manual prescribes.
\end{itemize}
\end{refusalbox}

Related: \menu{1}'s SCF-related findings (a report can tell you to come here),
and Chapter~\ref{ch:criteria} for the criteria the triage applies.
